% =====================================================================
% Referências (lista manual; a etapa 4 troca por biblatex-abnt + referencias.bib)
% =====================================================================
\begin{thebibliography}{99}
\addcontentsline{toc}{section}{Referências}

\bibitem[ANDERSON et al.(1999)]{anderson1999}
ANDERSON, E. et al. \textbf{LAPACK users' guide}. 3. ed. Philadelphia: SIAM, 1999.

\bibitem[BACKUS(1979)]{backus1979}
BACKUS, John. The history of FORTRAN I, II, and III. \textbf{IEEE Annals of the History of Computing}, v. 1, n. 1, p. 21--37, 1979.

\bibitem[BEZANSON et al.(2017)]{bezanson2017}
BEZANSON, Jeff; EDELMAN, Alan; KARPINSKI, Stefan; SHAH, Viral B. Julia: a fresh approach to numerical computing. \textbf{SIAM Review}, v. 59, n. 1, p. 65--98, 2017.

\bibitem[BOYCE; DIPRIMA(2012)]{boyce2012}
BOYCE, William E.; DIPRIMA, Richard C. \textbf{Elementary differential equations and boundary value problems}. 10. ed. Hoboken: Wiley, 2012.

\bibitem[BURDEN; FAIRES; BURDEN(2016)]{burden2016}
BURDEN, Richard L.; FAIRES, J. Douglas; BURDEN, Annette M. \textbf{Numerical analysis}. 10. ed. Boston: Cengage Learning, 2016.

\bibitem[BUTCHER(2016)]{butcher2016}
BUTCHER, John C. \textbf{Numerical methods for ordinary differential equations}. 3. ed. Chichester: John Wiley \& Sons, 2016.

\bibitem[BUTCHER; WANNER(1996)]{butcher1996}
BUTCHER, John C.; WANNER, Gerhard. Runge-Kutta methods: some historical notes. \textbf{Applied Numerical Mathematics}, v. 22, n. 1--3, p. 113--151, 1996.

\bibitem[CHAPRA; CANALE(2015)]{chapra2015}
CHAPRA, Steven C.; CANALE, Raymond P. \textbf{Numerical methods for engineers}. 7. ed. New York: McGraw-Hill Education, 2015.

\bibitem[COMPUTER HISTORY MUSEUM(s.d.)]{chm}
COMPUTER HISTORY MUSEUM. \textbf{A key component: the operational amplifier}. Mountain View. Disponível em: \url{https://www.computerhistory.org/revolution/analog-computers/3/154}. Acesso em: 23 nov. 2026.

\bibitem[EULER(1768)]{euler1768}
EULER, Leonhard. \textbf{Institutionum calculi integralis}. Petropolis: Academia Imperialis Scientiarum, 1768. v. 1.

\bibitem[FRANCO(2006)]{franco2006}
FRANCO, Neide Bertoldi. \textbf{Cálculo numérico}. São Paulo: Pearson Prentice Hall, 2006.

\bibitem[FRIGO; JOHNSON(2005)]{frigo2005}
FRIGO, Matteo; JOHNSON, Steven G. The design and implementation of FFTW3. \textbf{Proceedings of the IEEE}, v. 93, n. 2, p. 216--231, 2005.

\bibitem[GOLDBERG(1991)]{goldberg1991}
GOLDBERG, David. What every computer scientist should know about floating-point arithmetic. \textbf{ACM Computing Surveys}, v. 23, n. 1, p. 5--48, 1991.

\bibitem[GOLDSTINE(1972)]{goldstine1972}
GOLDSTINE, Herman H. \textbf{The computer from Pascal to von Neumann}. Princeton: Princeton University Press, 1972.

\bibitem[GOMEZ(1999)]{gomez1999}
GOMEZ, Claude (ed.). \textbf{Engineering and scientific computing with Scilab}. Boston: Birkhäuser, 1999.

\bibitem[HAIRER; NØRSETT; WANNER(1993)]{hairer1993}
HAIRER, Ernst; NØRSETT, Syvert P.; WANNER, Gerhard. \textbf{Solving ordinary differential equations I: nonstiff problems}. 2. ed. Berlin: Springer, 1993.

\bibitem[HARRIS et al.(2020)]{harris2020}
HARRIS, Charles R. et al. Array programming with NumPy. \textbf{Nature}, v. 585, p. 357--362, 2020.

\bibitem[HEUN(1900)]{heun1900}
HEUN, Karl. Neue Methode zur approximativen Integration der Differentialgleichungen einer unabhängigen Veränderlichen. \textbf{Zeitschrift für Mathematik und Physik}, v. 45, p. 23--38, 1900.

\bibitem[HIGHAM(2002)]{higham2002}
HIGHAM, Nicholas J. \textbf{Accuracy and stability of numerical algorithms}. 2. ed. Philadelphia: SIAM, 2002.

\bibitem[HOPP(1999)]{hopp1999}
HOPP, Peter M. \textbf{Slide rules: their history, models and makers}. Mendham: Astragal Press, 1999.

\bibitem[IEEE(2019)]{ieee2019}
IEEE. \textbf{IEEE Std 754-2019}: IEEE standard for floating-point arithmetic. New York: IEEE, 2019.

\bibitem[JUNG(2005)]{jung2005}
JUNG, Walter G. (ed.). \textbf{Op amp applications handbook}. Burlington: Newnes, 2005.

\bibitem[KUTTA(1901)]{kutta1901}
KUTTA, Wilhelm. Beitrag zur näherungsweisen Integration totaler Differentialgleichungen. \textbf{Zeitschrift für Mathematik und Physik}, v. 46, p. 435--453, 1901.

\bibitem[LATHI; DING(2009)]{lathi2009}
LATHI, B. P.; DING, Zhi. \textbf{Modern digital and analog communication systems}. 4. ed. New York: Oxford University Press, 2009.

\bibitem[LAWSON; HANSON; KINCAID; KROGH(1979)]{lawson1979}
LAWSON, Charles L.; HANSON, Richard J.; KINCAID, David R.; KROGH, Fred T. Basic linear algebra subprograms for Fortran usage. \textbf{ACM Transactions on Mathematical Software}, v. 5, n. 3, p. 308--323, 1979.

\bibitem[LEIBNIZ(1703)]{leibniz1703}
LEIBNIZ, Gottfried Wilhelm. Explication de l'arithmétique binaire. \textbf{Mémoires de mathématique et de physique de l'Académie royale des sciences}, Paris, 1703.

\bibitem[LENSTRA; LENSTRA JR.; LOVÁSZ(1982)]{lenstra1982}
LENSTRA, A. K.; LENSTRA JR., H. W.; LOVÁSZ, L. Factoring polynomials with rational coefficients. \textbf{Mathematische Annalen}, v. 261, p. 515--534, 1982.

\bibitem[L'ECUYER(2017)]{lecuyer2017}
L'ECUYER, Pierre. History of uniform random number generation. In: WINTER SIMULATION CONFERENCE, 2017, Las Vegas. \textbf{Proceedings} [...]. Piscataway: IEEE, 2017.

\bibitem[MENABREA(1843)]{menabrea1843}
MENABREA, Luigi Federico. Sketch of the Analytical Engine invented by Charles Babbage. Tradução e notas de Augusta Ada King, Condessa de Lovelace. \textbf{Scientific Memoirs}, Londres, v. 3, p. 666--731, 1843.

\bibitem[MOLER(2004)]{moler2004}
MOLER, Cleve B. \textbf{Numerical computing with MATLAB}. Philadelphia: SIAM, 2004.

\bibitem[NAGEL(1975)]{nagel1975}
NAGEL, Laurence W. \textbf{SPICE2: a computer program to simulate semiconductor circuits}. Berkeley: University of California, Electronics Research Laboratory, 1975. Memorandum n. ERL-M520.

\bibitem[NILSSON; RIEDEL(2019)]{nilsson2019}
NILSSON, James W.; RIEDEL, Susan A. \textbf{Electric circuits}. 11. ed. Boston: Pearson, 2019.

\bibitem[OPPENHEIM; WILLSKY; NAWAB(1997)]{oppenheim1997}
OPPENHEIM, Alan V.; WILLSKY, Alan S.; NAWAB, S. Hamid. \textbf{Signals and systems}. 2. ed. Upper Saddle River: Prentice Hall, 1997.

\bibitem[OVERTON(2001)]{overton2001}
OVERTON, Michael L. \textbf{Numerical computing with IEEE floating point arithmetic}. Philadelphia: SIAM, 2001.

\bibitem[PRESS; TEUKOLSKY; VETTERLING; FLANNERY(2007)]{press2007}
PRESS, William H.; TEUKOLSKY, Saul A.; VETTERLING, William T.; FLANNERY, Brian P. \textbf{Numerical recipes: the art of scientific computing}. 3. ed. Cambridge: Cambridge University Press, 2007.

\bibitem[QUARTERONI; SACCO; SALERI(2007)]{quarteroni2007}
QUARTERONI, Alfio; SACCO, Riccardo; SALERI, Fausto. \textbf{Numerical mathematics}. 2. ed. Berlin: Springer, 2007.

\bibitem[RAND CORPORATION(1955)]{rand1955}
RAND CORPORATION. \textbf{A million random digits with 100,000 normal deviates}. Glencoe: Free Press, 1955.

\bibitem[ROJAS(1997)]{rojas1997}
ROJAS, Raúl. Konrad Zuse's legacy: the architecture of the Z1 and Z3. \textbf{IEEE Annals of the History of Computing}, v. 19, n. 2, p. 5--16, 1997.

\bibitem[RUGGIERO; LOPES(1996)]{ruggiero1996}
RUGGIERO, Márcia A. Gomes; LOPES, Vera Lúcia da Rocha. \textbf{Cálculo numérico: aspectos teóricos e computacionais}. 2. ed. São Paulo: Makron Books, 1996.

\bibitem[RUNGE(1895)]{runge1895}
RUNGE, Carl. Über die numerische Auflösung von Differentialgleichungen. \textbf{Mathematische Annalen}, v. 46, p. 167--178, 1895.

\bibitem[SCILAB(2013)]{scilab2013}
SCILAB. \textbf{History}. 2013. Disponível em: \href{https://web.archive.org/web/20130317145202/http://www.scilab.org/scilab/history}{\nolinkurl{https://web.archive.org/web/20130317145202/}\penalty-10000 \nolinkurl{http://www.scilab.org/scilab/history}}. Acesso em: 23 nov. 2026.

\bibitem[SWADE(2001)]{swade2001}
SWADE, Doron. \textbf{The difference engine: Charles Babbage and the quest to build the first computer}. New York: Viking, 2001.

\bibitem[VIRTANEN et al.(2020)]{virtanen2020}
VIRTANEN, Pauli et al. SciPy 1.0: fundamental algorithms for scientific computing in Python. \textbf{Nature Methods}, v. 17, p. 261--272, 2020.

\bibitem[VON NEUMANN(1945)]{vonneumann1945}
VON NEUMANN, John. \textbf{First draft of a report on the EDVAC}. Philadelphia: Moore School of Electrical Engineering, University of Pennsylvania, 1945.

\bibitem[WILLIAMS(1997)]{williams1997}
WILLIAMS, Michael R. \textbf{A history of computing technology}. 2. ed. Los Alamitos: IEEE Computer Society Press, 1997.

\end{thebibliography}
